\section{Disclosed task, mechanism and delivered information}
\label{app:task}
Represent the two-dimensional state and action by complex numbers $z_t$ and $u_t$. A step is indexed by its pre-action state: $u_0$ produces $z_1$. The initial conditions are $z_0=1$, $\phi_0=4.147589101011308$ and $u_{-1}=0$. The evaluated world obeys
\begin{align}
 d(z)&=(1.06+0.015|z|^2)z+(0.12-0.08i),\\
 z_{t+1}&=d(z_t)+e^{i\phi_t}u_t,\qquad |u_t|\leq 0.6,\\
 \phi_{t+1}&=\phi_t+\omega+\kappa\,\operatorname{Re}(u_{t-1})\operatorname{Im}(u_{t-1}),\qquad\omega=0.1.
\end{align}
The lag in the phase update is part of the mechanism: the update uses $u_{t-1}$, not the action just chosen. The original family allowed $\kappa\in\{0.3,3,6,12,18\}$. Only the two settings below occur in the fixed-history pilot. They are parameter settings, not a validated difficulty ordering.

\begin{center}
\begin{tabular}{llll}
\toprule
Paper alias & Archival case identifier & Snapshot after step & $\kappa$\\
\midrule
A & case-01 & 4 & 12\\
B & boundary-01 & 11 & 12\\
C & case-03 & 4 & 18\\
D & boundary-02 & 9 & 18\\
\bottomrule
\end{tabular}
\end{center}

An episode ends on escape ($|z_t|>4$) or after 32 actions. Success requires completion and $|z_t|\leq0.05+10^{-10}$ at \emph{all} $t\in\{28,29,30,31,32\}$. Merely avoiding escape is insufficient. Action validation rejects nonfinite coordinates or norms above $0.6+10^{-9}$; a smaller tolerance overshoot is normalized to $0.6$. State observations are rounded to twelve decimal places, whereas the simulator retains binary floating-point state. Stored actions retain their applied precision. The exact numeric history of every selected case is included in \texttt{legacy-replay/data/cases.json}; it is not regenerated from a caption or rounded manuscript table.

The participant receives the objective, horizon, action limit, escape threshold, full permitted action--state history and remaining analysis allowance. It does not receive these equations, $\phi_t$, $\omega$, $\kappa$, the evaluator's continuation or its loss diagnosis. The opaque identifiers in delivered inputs do not disclose the boundary/comparison labels. The frozen selector and schemas are inspectable in the reader artifact; the full provider logs are not redistributed. Disclosure to reviewers does not alter the historical participant information boundary.

\section{Witness and loss-bound justification}
\label{app:bounds}
\paragraph{Positive witness.}
For every pre-action snapshot and every resulting state, one fixed privileged controller attempts a continuation. At each step it chooses
\begin{equation}
 u=\operatorname{clip}_{0.6}\!\left[-d(z)e^{-i\phi}\right],\qquad
 \operatorname{clip}_{r}(v)=
 \begin{cases}
   v, & \lVert v\rVert\leq r,\\[3pt]
   \displaystyle\frac{r v}{\lVert v\rVert}, & \lVert v\rVert>r.
 \end{cases}
\end{equation}
Here $\lVert v\rVert$ is the Euclidean norm, equal to the modulus for a complex scalar.
The original transition, action validation, escape check and scheduled target test score the entire continuation. A successful admissible continuation establishes existence for the implemented task. A failed continuation does not establish impossibility. This controller knows the hidden phase and mechanism; its success is not evidence of blind action identification or participant learning. Every pilot pre-action state passed this positive baseline.

\paragraph{Sufficient negative bound for the ideal recurrence.}
Write \(x = (\operatorname{Re} z, \operatorname{Im} z)\) and $c=(0.12,-0.08)$. Fix a vector $q$ with $\|q\|_2\leq1$. Suppose $q^Tx\geq p>0$. Since $\|x\|_2\geq q^Tx\geq p$, and the scalar multiplier $1.06+0.015\|x\|_2^2$ is positive, the next-state projection satisfies
\begin{align}
q^Tx'&=(1.06+0.015\|x\|_2^2)q^Tx+q^Tc+q^TR(\phi)u\\
&\geq (1.06+0.015p^2)p+q^Tc-0.6.
\end{align}
The last term follows from Cauchy--Schwarz and the admissible action radius. The function $g(p)=(1.06+0.015p^2)p$ is increasing on the positive domain. Consequently, iterating this lower bound while it remains positive gives a bound for \emph{every} allowed future control sequence, irrespective of the future phases. A bound exceeding 4 before the horizon implies unavoidable escape; a bound exceeding the target tolerance at a required target step excludes successful completion. A target violation already present in the recorded prefix also suffices, because future actions cannot change a past required state.

\paragraph{Implementation and numerical scope.}
The evaluator chooses $q$ by truncating each coordinate of $x/\|x\|$ toward zero to four decimal places, and checks its norm. It initializes $p_0=q^Tx$ from the stored full-precision state. The implemented recurrence is
\begin{equation}
p_{j+1}=\left\lfloor (1.06+0.015p_j^2)p_j+q^Tc-0.6-10^{-9}\right\rfloor_{10^{-12}},
\end{equation}
where $\lfloor\cdot\rfloor_{10^{-12}}$ rounds downward to that grid. Coefficients and initial state are converted from their actual binary-float values to exact rational numbers. Loss is reported only for a positive induction sequence crossing $4.000001$ within the remaining horizon, or crossing $0.05+10^{-10}$ at a required target step. A separate 110-digit decimal implementation checks the direction, every bound, the positivity domain, threshold and horizon. The code also checks that the transition definition and task constants match the bound's assumptions.

\paragraph{Archived arithmetic audit.}
A subsequent record-only audit covers the pilot, distinct component
grades and the continuous subset of the prospective panel. All 50 distinct
stored directions, including directions for which no loss was certified,
have exact squared norm at most one; the minimum value of
$1-\|q\|_2^2$ is $5.75\times10^{-6}$. There are ten distinct certified
projection records (eleven occurrences). Rational recomputation reproduces
the archived binary64 bounds exactly. Recomputing these recurrences with per-step
subtractions $10^{-9}$, $2\times10^{-9}$ and $10^{-8}$ preserves every
certificate kind and first certification step. Even at the enlarged
slacks, the smallest final lower-bound gap above the applicable runtime
threshold exceeds $0.0371$. No simulator transition or model request is
performed by this audit, and no historical certificate is rewritten.

The action validator rejects computed norms above $0.6+10^{-9}$ and
rescales a smaller computed overshoot by $a\,0.6/|a|$. It therefore does
not simply apply every action up to the rejection tolerance unchanged.
For the 1,201 distinct applied actions in the inspected continuous
trajectory records, exact squared-norm checks give a maximum norm of
approximately $0.60000000000000002694$. This finite record check does not
bound all possible future actions or transcendental-function errors.

\paragraph{Runtime-specific conditional budget.}
A separate local audit inspects the pinned interpreter and the currently
loaded ARM64 math image. Its compiled complex multiplication uses fused
multiply-add instructions; its complex exponential calls
\texttt{exp} and \texttt{\_\_sincos\_stret}. The norm uses platform
\texttt{hypot}, and the square still uses platform \texttt{pow}.
Thus a source-level list of separate sine/cosine calls does not fully
describe this executable.

Appendix~\ref{app:runtime-proof} derives a norm bound for the captured
\texttt{hypot} instruction sequence under ordinary IEEE binary64
round-to-nearest semantics with gradual underflow. It then reduces the
remaining library accuracy requirements to an absolute square error
$\epsilon_{\rm pow}$ and a norm excess $\delta_{\rm rot}$ for the
computed complex rotation. Accuracy relative to the true angle is not
needed for this projection bound. Taking each of these two bounds as
$2^{-30}$ yields total one-step error below $6.288\times10^{-10}$,
which fits the archived $10^{-9}$ subtraction. The local checker verifies
the scalar inequalities with exact rational arithmetic.

This is a sufficient conditional bound, not an attestation that the two
library bounds hold uniformly. Finite current-process probes satisfy them
but do not prove them. In addition, the two inspected historical runtime
inventories and activation records do not pin the loaded math-image
identity or floating-point control register. Today's image capture cannot
retroactively fill that gap. The ten archived projection certificates
therefore retain their stated conditional numerical status. The route
family's exact integer transitions and finite-horizon dynamic program do
not use this projection premise.

The classifier returns preserved only with a successful witness, lost only with a sufficient negative certificate, and unresolved otherwise. Simultaneous positive and negative evidence is treated as an error. This is an incomplete classifier, not a complete viability kernel or a map of a continuous safe cone. A numerical failure in candidate optimization is a different unresolved category and is not fed to this classifier as physical loss.

\section{Computational interventions and reproducible selectors}
\label{app:components}
The candidate coefficients and original analysis outputs are fixed data. No candidate is refitted after a component is crossed. Let $f_m+B_m u$ denote candidate $m$'s predicted next state at the selected snapshot. Candidate construction, feature families, phase extrapolation and the exact numeric coefficients are inspectable in \texttt{legacy-replay/data/component\_inputs.json}. The replay implements these fixed operations; it does not execute the archived participant programs.

\paragraph{One-step worst-case rule.}
For $G$ equally spaced angles and radius $r$, set $u_j=r(\cos(2\pi j/G),\sin(2\pi j/G))$ and choose the first grid index minimizing $\max_m\|f_m+B_m u_j\|_2$. Snapshot B uses $G = 720$ and $r = 0.6$; D uses $G = 1440$ and $r = 0.599999999$. The four finer-grid checks use $G=2880$. This is optimization on a finite circle, not a guarantee of global optimization over the full action disk or of task completion.

\paragraph{Past-error weighted rule.}
Each fitted parent has weight $w_m=1/(10^{-4}+e_m^2)$, where $e_m$ is its recorded leave-one-out residual RMS. In the expanded D component, six scenarios inherit equal shares $w_m/6$ of each parent's weight. With $v=\sum_m w_m B_m^Tf_m$, the selected action is $u=-0.6v/\|v\|$. The matrices here are scalar gains times rotations, so the quadratic term in $\sum_mw_m\|f_m+B_mu\|^2$ is constant on the radius-$0.6$ circle. Thus this direction minimizes that fixed-radius surrogate for these matrices. It is not a general theorem for arbitrary linear-control matrices, a calibrated posterior, or a worst-case task-safety guarantee. D's second representation comprises 11 parent fits expanded by two phase modes and three phase offsets $\{-0.12,\,0,\,0.12\}$; the complete construction is in the code and numeric inputs.

\paragraph{Multi-step plans and the researcher completion.}
For each of B's candidates, the individual planner optimizes the remaining $N=32-n$ controls with projection to the action disk. Its surrogate objective is
\begin{equation}
 J_m=\sum_{i=1}^{N}\left[w_{n+i}\|\hat x_{n+i}\|^2+\lambda\big(\max\{0,\|\hat x_{n+i}\|^2-14.44\}\big)^2+10^{-5}\|u_{n+i-1}\|^2\right],
\end{equation}
where $w_t=1$ for $t\geq28$ and $10^{-4}$ otherwise, and $\lambda=10$. The deterministic initialization, projected accelerated-gradient updates, backtracking limits and restart schedule are included in \texttt{components.py}; individual plans use 1,000 iterations. The declared selector R2* chooses the individual plan with the lowest computed cost and takes its first action. This final selection rule was completed by the researchers. Agreement with the original choice is a reconstruction check, not proof that it was the model's latent policy. At least one crossed action has magnitude approximately $0.5755$, so this contrast does not isolate direction or horizon while holding magnitude fixed.

The two auxiliary joint plans use a shared first action, separate candidate-specific later actions, the mean of the candidate costs, and 1,400 iterations. One is numerically unresolved. The one-at-a-time D deletions remove each of 66 scenarios or each of 11 six-scenario parent groups; all other settings stay fixed and no refitting occurs. No-rescue results cover only these interventions, not joint deletion, changed fitting or arbitrary alternative uncertainty models.

The 91 planned cells comprise 8 primary crossings, 4 finer-grid checks, 2 joint-plan auxiliaries and 77 omissions. Ninety were computed and one auxiliary remained numerically unresolved. Repeated actions are deduplicated before privileged grading, giving 13 distinct grades. There is no additional model sample in this count. The full records, rather than the rounded values in figures, define action identity.

\subsection{Five retrospective completions at B}
\label{app:b-completions}
This separate sensitivity analysis holds the eight individual plans from
B's native-loss pool A1 and four from its native-preservation pool A2 fixed.
It changes only which plan
supplies the first action. For a plan's predicted states $\hat x_t$, the
five minimized criteria are: the total penalized cost $J_m$ above;
$\max_{28\leq t\leq32}\|\hat x_t\|$;
$\tfrac15\sum_{t=28}^{32}\|\hat x_t\|^2$;
$\max_{n\leq t\leq32}\|\hat x_t\|$, including the common starting state;
and $\|\hat x_{32}\|$. All ties use the lowest existing plan index.
The last criterion covers only the final target step; it is not the full
task-success condition. The finite selector list was fixed before its
additional grading, with the original outcomes already known.

\begin{table}[htbp]
\centering\small
\begin{tabular}{lccr}\toprule
Criterion minimized & Native A2 & Crossed A1 & A1 action norm\\\midrule
Total penalized cost (original R2*) & P & P & 0.575465\\
Maximum target-step norm & P & P & 0.575465\\
Mean squared target-step norm & P & P & 0.575465\\
Trajectory peak norm & P & L & 0.042684\\
Final-step norm & P & P & 0.579114\\\bottomrule
\end{tabular}
\caption{B selector sensitivity from the archived fixed-plan analysis.
All five reproduce the native A2 action within $10^{-6}$; the first three
choose the same A1 plan. The minimum-peak completion reverses the reported
rescue. Selector counts are not probabilities or independent replications.}
\label{tab:b-completions}
\end{table}

Thus native-action reproduction alone does not identify the crossed
completion. The reversal also changes the chosen plan and action magnitude;
it does not isolate a horizon-only or direction-only effect. The saved
selector indices, exact actions, metrics and labels are in
\texttt{artifact/source/pilot-sensitivity.private.json}. This finite
sensitivity result does not cover new optimizers, mixtures or arbitrary
alternative policies.

\section{Protocol, record boundaries and cost accounting}
\label{app:protocol}
The fixed trial order was A1, D1, C1, B1, C2, B2, A2, D2, where the suffix denotes repeat. Each trial begins in a fresh local context with an opaque snapshot identifier; both repeats of a case receive exactly the same initial payload bytes. Each subsequent payload contains the same physical history and that trial's own submitted analysis code and returned results. Therefore the third or fourth requests are not equal-input repetitions of the first request.

The wire schema permits one analysis request or one final two-coordinate action. Analysis code must define \texttt{act(history)} and can use the allowed mathematical facilities; no network, filesystem, process creation or hidden simulator interface is provided. The protocol and policy source state the exact permitted syntax and names. Local analysis has 30-second CPU/wall limits, a sampled 1 GiB RSS stop and a 100,000-character output limit. Sampled RSS is not a hard allocation ceiling. Each trial permits at most three analysis requests and one final action, for at most four model requests in total. Trial wall time is 3,780 seconds; request time is bounded by the smaller of 900 seconds and remaining trial time minus a 120-second record/grading margin. No retry, neutral request, replacement or outcome-dependent stopping is permitted.

The model identifier and effort are configuration records, not attested server weights. Local receipts record one agent message per call, no unexpected tool events and cleanup of initially empty call directories. These checks do not establish independent sampling seeds or provider-side isolation. Full provider transcripts remain in the source archive; this bundle distributes selected numerical inputs and final record projections. Hidden chain-of-thought is not observed or claimed as evidence.

The recorded 588,051 input and 129,677 output tokens total 717,728. The 92,160 cached-input tokens are already included in input. Reasoning-output tokens are likewise a subset of output. There were 28 receipted requests and zero unreceipted requests in this pilot. Provider receipts are the source of these counts; analyst use and money are unmeasured. The historical technical null's unresolved reservation is not a zero-use observation and is not included in the pilot total.

\section{Prospective tasks, reference histories and adapter}
\label{app:transfer}
Exact configurations and twelve packets are in
\texttt{legacy-replay/data/transfer/}; the numeric replay inputs under
\texttt{legacy-replay/data/transfer/} are checked against these records.

\paragraph{Continuous reference histories.}
The recurrence and task target are those in Appendix A, with three new
initial phase/frequency configurations and couplings 3, 6 and 12.
Phases and frequencies derive deterministically from SHA-256 of the
literal seed \texttt{transfer-validation-20261025-v1} and instance
ordinal. A prescribed privileged ten-step path tracks waypoints with
radii $1.1,1.2,1.3,1.4,1.5,1.6,1.75,1.9,2.05,2.2$ and deterministic
changing angles. Snapshots are after steps 6 and 10. Pre-collection
checks require legal histories and a positive continuation, without
resampling. This selection does not establish calibrated boundary
distances or natural participant behavior.

\paragraph{Route--resource dynamics and certificates.}
The state is node $n\in\{0,\ldots,18\}$ and remaining integer fuel $f$.
Action $j\in\{0,1,2\}$ is encoded as $[j,0]$ and advances to
$n'=\min(18,n+j+1)$. The hidden stationary positive edge cost is
\begin{equation}
c(n,j)=b_j+((n+o_j)\bmod p_j),\qquad f'=f-c(n,j).
\end{equation}
Success means reaching node 18 within twelve actions without negative
fuel. The three fixed $(b,o,p)$ configurations and initial fuel values
are in \texttt{CASES.private.json}. Unlike continuous control, every
legal route action is in the menu. The participant knows the graph,
objective, horizon and observed costs through its action--state history,
but not the edge-cost formula or private coefficients.

Let $V_h(n)$ be the minimum fuel needed to reach the goal in at most
$h$ further actions. Set $V_0(18)=0$, $V_0(n<18)=+\infty$, and
\begin{equation}
V_h(18)=0,\qquad
V_h(n)=\min_{j\in\{0,1,2\}}
\{c(n,j)+V_{h-1}(\min(18,n+j+1))\}.
\end{equation}
The code represents infinity by \texttt{None}. With positive costs,
$f\geq V_h(n)$ provides a minimum-cost admissible continuation;
$f<V_h(n)$, negative fuel or an unreachable goal certifies loss.
Every value-table cell is checked by a separately implemented recursive
recurrence. Positive witnesses are replayed. Initial fuel is the exact
minimum completion cost plus one; the lexicographically first
minimum-cost path supplies snapshots after steps 2 and 4. Baseline
feasibility does not prove that the hidden cost law is identifiable from
the observed short prefix.

\paragraph{Exact optional table contract.}
The diagnostic continuous menu has twelve equally spaced directions at
radius $0.6$, rounded to twelve decimals, plus zero. For a returned
table with losses $L_{mj}$ and weights $w_m$, the weighted rule chooses
the first menu index minimizing $\sum_m w_mL_{mj}/\sum_mw_m$; the
worst-case rule minimizes $\max_m L_{mj}$ with the same tie break.
Weights belong to the candidate pool. The adapter accepts only the
exact keys \texttt{actions}, \texttt{candidate\_losses},
\texttt{weights}, and \texttt{rule}, the exact delivered menu, 1--128
candidate rows and finite entries of absolute magnitude at most
$10^{100}$. Weights are finite, nonnegative and bounded by $10^{100}$,
with positive total. Final-action reproduction uses maximum coordinate
error $\leq10^{-8}$. The last successful analysis is selected before
examining outcome category. A successful earlier table cannot rescue
a missing or unfaithful final table.

The replay changes only imports/data paths, preserves frozen numeric
functions, and omits dispatch code. Original function text and
source/export digests are included. It reads bound archived analysis
values; participant programs remain unexecuted text. It reproduces the
eligibility test without inferring or regenerating tables.
\section{Narrow runtime proof audit}
\label{app:runtime-proof}

This appendix states a sufficient numerical contract for the continuous
projection argument and documents which parts were established locally.
The real-arithmetic argument, the archived recurrence computation, the
current machine implementation and the historical execution environment
are distinct evidence layers. None of the checks below adds a live agent
trial or reruns a physical episode.

\paragraph{Identity and trust assumptions.}
The inspected executable matches the experimental Python digest. It is
CPython 3.12.14 on ARM64. The current math image is
\texttt{libsystem\_m.dylib}; its UUID, text/constants hashes and relevant
function bytes were recorded. The current thread reports a zero floating-point
control register: round to nearest, with no flush-to-zero or alternative
handling enabled. We assume the ordinary IEEE binary64 semantics of
the decoded add, multiply, divide, fused multiply-add and square-root
instructions. The source-to-instruction interpretation is a manual
derivation with exact constant checks, not a proof-assistant-certified
binary analysis. The historical inventories pin Python but do not attest
the loaded math image or control register; equivalence to the historical
environment remains an explicit assumption.

\paragraph{Bound for the captured norm implementation.}
Let $u=2^{-53}$ and $d=2^{-1075}$, the unit roundoff and half the smallest
subnormal. Here $u$ denotes roundoff, not a physical action. For finite
binary64 coordinates, put $a=\max(|x_1|,|x_2|)$ and
$b=\min(|x_1|,|x_2|)$, and let $r=\sqrt{a^2+b^2}$ in exact arithmetic.
The captured function returns $a$ when $b=0$. Otherwise it drops $b$
only if the stored exponent difference exceeds 27. This implies
$b/a<2^{-27}$, including a subnormal $b$, so dropping it reduces the
exact norm by less than $r2^{-55}=ru/4$.

For nonzero retained inputs, the decoded branches use the following
scales, where $e_a$ is the biased exponent field of $a$:
\begin{center}
\begin{tabular}{@{}lll@{}}
\toprule
Exponent field & Input scale & Undo scale\\
\midrule
$e_a<539$ & $2^{768}$ & $2^{-768}$\\
$539\leq e_a<1534$ & $1$ & $1$\\
$e_a\geq1534$ & $2^{-768}$ & $2^{768}$\\
\bottomrule
\end{tabular}
\end{center}
The retained inputs scale exactly. In the small-input branch even the
smallest subnormal becomes $2^{-306}$. In the large-input branch a
nonzero retained $b$ has exponent difference at most 27 and remains
normal after scaling. In all branches the squared larger input and
the sum are normal and cannot overflow. The core computes
$t=\operatorname{RN}(a_s^2)$,
$v=\operatorname{RN}(b_s^2+t)$ by a fused operation, and
$w=\operatorname{RN}(\sqrt v)$, before undoing the scale.
Thus $w$ lies within factors $(1-u)^2$ and $(1+u)^2$ of the scaled norm.
Undoing a power-of-two scale is exact for a finite normal output;
subnormal rounding contributes at most $d$. Overflowed outputs are not
finite accepted norms. Combining the discarded-input allowance gives
the sufficient bound, for a finite returned norm $h$,
\begin{equation}
 r(1-3u)-d\ \leq\ h\ \leq\ r(1+3u)+d.
 \label{eq:hypot-bound}
\end{equation}
The rational checker verifies
$(1-u/4)(1-u)^2\geq1-3u$ and $(1+u)^2\leq1+3u$.
In particular, $h\leq4$ implies $r<5$ and
$|h-r|<H=16u$. This is an analytic bound for the captured sequence
under the stated instruction semantics, not an accuracy inference from
test samples.

\paragraph{Action normalization.}
Let $L$ be the exact binary64 value of $0.6$. An accepted but unscaled
action satisfies $h\leq L$, hence
$r\leq(L+d)/(1-3u)<L+8u$. When $h>L$, the validator rescales by
$aL/h$. Here $r>1/2$ and $h>1/2$. The compiled multiplication by a real
scalar and division by a positive real denominator each round each
component once; the zero imaginary terms do not introduce a separate
nonzero error. Using $\sqrt2<3/2$ gives
\begin{equation}
 \|a_{\rm applied}\|_2
 \leq \frac{L(1+u)^2}{1-3u-2d}
       +3d(1+u)+\tfrac32 d
 < L+8u.
 \label{eq:action-bound}
\end{equation}
This covers accepted rescaling, including subnormal allowances, rather
than treating the rejection tolerance as an applied-action allowance.

\paragraph{Two remaining library accuracy contracts.}
Write $\widehat s=\operatorname{pow}(h,2)$ and
$C(\phi)=\operatorname{cmath.exp}(i\phi)$ for the computed values.
The needed contracts are
\begin{equation}
 |\widehat s-h^2|\leq\epsilon_{\rm pow},\qquad
 \|C(\phi)\|_2\leq1+\delta_{\rm rot}.
 \label{eq:remaining-contracts}
\end{equation}
They must hold uniformly over relevant norm inputs and permitted finite
phases. The second includes the wrapper's \texttt{exp} and combined
\texttt{\_\_sincos\_stret} computation. No bound on angular error is
needed: every such control vector has norm at most
$(1+\delta_{\rm rot})(L+8u)$, irrespective of its direction.
Inspection of call targets establishes the operation path but does not
prove either accuracy contract.

\paragraph{Composed one-step bound.}
For a state before computed escape, $r<5$. Let $\alpha,\beta,c,L$
denote the exact binary64 task constants, so the reference drift is
$(\alpha+\beta r^2)x+c$. Set $B=32u+d$. Under the sufficient choices
$\epsilon_{\rm pow}=\delta_{\rm rot}=2^{-30}$, the relevant basic-operation
intermediates have magnitude below 32; the square is handled separately
by its library contract. Each basic rounding error is therefore at most
$B$. The squared-radius error is bounded by
$10H+H^2+\epsilon_{\rm pow}$. Two basic operations form the drift
multiplier. Scaling the state and adding the bias give vector error
at most $5\beta(10H+H^2+\epsilon_{\rm pow})+13B$.
The compiled complex product uses one ordinary product and one fused
operation per component; its error contributes less than $3B$.
Final vector addition contributes less than $3B/2$.

Consequently the sufficient loss allowance beyond the ideal $-L$ term is
\begin{equation}
 E=5\beta(10H+H^2+\epsilon_{\rm pow})
    +18B+8u+\delta_{\rm rot}(L+8u).
 \label{eq:composed-error}
\end{equation}
The exact rational evaluation is approximately
$6.2870891\times10^{-10}<10^{-9}$. For $\|q\|_2\leq1$ and a positive
projection lower bound $p$, this yields
\[
 q^\top x'\geq(\alpha+\beta p^2)p+q^\top c-L-E.
\]
The archived subtraction and downward floor therefore remain sufficient
\emph{if} the stated instruction, library and historical-equivalence
premises hold. The prior audit's minimum terminal gap exceeds $0.0371$;
it also clears the norm-test error. The same ten unique archived
projection certificates are covered conditionally, with no relabeling.

\paragraph{Checks and open obligations.}
The archived local audit checked exact inequality constants, captured function
identity and historical inventory coverage. The reader bundle rechecks only
the exact constants and exponent ranges with \texttt{artifact/numeric\_constants.py};
it does not recapture or attest the historical runtime. Fifty-two norm inputs,
including exponent/underflow edges, twenty square inputs and twenty-six
phases were probed in the archived audit interpreter. Every finite tested output
satisfies the sufficient bounds; an overflowed norm is recorded separately.
These are scalar diagnostics, not uniform accuracy proofs. Negative
controls reject altered function bytes and oversized error budgets, and
retain an open status when a library premise or historical binding is
missing. They test the checker, not the missing premises themselves.

The remaining obligations are explicit: prove both bounds in
Eq.~\eqref{eq:remaining-contracts} for the relevant library implementation,
and justify its equivalence and floating-point settings in the historical
runs. Current image hashes cannot supply absent historical attestation.
Accordingly, this addendum narrows the numerical assumptions and supports
their composition; it does not turn the archived results into unconditional
machine-verified claims. Exact route dynamic programming and the
observed trial counts are unchanged.


\section{The 32-trial common-menu protocol and panel usage}
\label{app:common-menu-protocol}

\paragraph{Selected histories and menus.}
Table~\ref{tab:common-menu-histories} identifies the four purposeful
histories using the earlier transfer panel's source IDs. The two
continuous histories are at physical step 10; the two route histories
are at step 2. Each appears with both profiles, both repeats and both
conditions. These are selected reference histories, not a random sample
of histories or eight newly sampled worlds.

\begin{table}[htbp]
\centering\small
\begin{tabular}{lllrrr}\toprule
Delivered ID & Source ID & Family & Instance & Cutoff & Menu size\\\midrule
probe-01 & snapshot-03 & Continuous & 1 & 10 & 13\\
probe-02 & snapshot-07 & Continuous & 2 & 10 & 13\\
probe-03 & snapshot-02 & Route & 1 & 2 & 3\\
probe-04 & snapshot-06 & Route & 2 & 2 & 3\\\bottomrule
\end{tabular}
\caption{Fixed histories for the 32-trial panel. Source histories and
menu identities are hash-bound in the archived allocation.}
\label{tab:common-menu-histories}
\end{table}

Both conditions receive identical history and task/menu bytes for a
given probe; the condition instruction differs. The continuous menu
is the twelve radius-$0.6$ directions at angles $2\pi j/12$,
$j=0,\ldots,11$, rounded to twelve decimals, followed by zero.
The route menu is $[(0,0),(1,0),(2,0)]$ in that order. Only the indexed
menu action is executed. The original physical clock, remaining
horizon and task objective are unchanged. The participant receives
visible history, public task/menu and its own analysis outputs, without
hidden dynamics or future observations. Trial order was fixed before
the run and is distributed in \texttt{artifact/data/crossing.json}.

\paragraph{Common output scaffold and condition-specific execution.}
The shared structured-output scaffold wraps a \texttt{decision} object
that can request analysis or return an action ID, a record, or both.
It does not itself enforce condition-specific fields, record identity,
dimensions, loss range or action agreement; the local contract does.
A record has exactly six fields: \texttt{schema}, \texttt{case\_id},
\texttt{menu\_sha256}, \texttt{loss\_definition\_id},
\texttt{candidate\_losses} and \texttt{rule}. It contains 1--128 rows,
one numeric loss in $[0,1]$ per menu action, and no weights. Each row
uses the same declared meaning: an uncalibrated forecast of failing the
original full task after that next action and the best continuation the
participant can plan. The two rules minimize the uniform row mean or
row maximum. Decimal numeric tokens are evaluated as exact rationals;
ties select the lowest menu index. Single and duplicate rows are legal,
and duplicates retain multiplicity. Rows are not verified world models.

In the optional condition, a legal \texttt{action\_id} is required and
an accompanying record is optional. Record failure or disagreement is
logged without replacing that action. In the mandatory condition, only
the record is allowed; its declared rule deterministically selects the
executed action. An invalid record produces no action and no retry.
Thus mandatory agreement is by construction, while optional agreement
is an observational check. This six-field contract differs from the
48-trial panel's weighted, four-field last-analysis adapter in Appendix E;
their eligibility counts are not measurements of the same interface.

\paragraph{Budgets and recorded use.}
For the 32-trial panel, each fresh trial permits up to three optional
local analyses and one final decision: at most four model requests,
3,780 seconds per trial and 900 seconds per request. Analysis has a
30-second CPU/wall limit, sampled 1 GiB RSS stop and 100,000-character
JSON output cap. It yields no new physical observation. No retry,
repair prompt, neutral request or outcome-dependent extension is
allowed; the panel ceiling is 128 model requests. The completed panel
used 48 analysis requests and 32 final decisions. Its single selected
action per trial was graded; crossed labels reuse the bound menu census,
not additional live observations.

\begin{table}[htbp]
\centering\small
\begin{tabular}{lrrrrr}\toprule
Panel & Requests & Input & Output & Cached input & Unknown calls\\\midrule
48-trial transfer & 120 & 2,325,280 & 484,134 & 372,352 & 0\\
32-trial common menu & 80 & 1,733,978 & 422,952 & 263,168 & 0\\\bottomrule
\end{tabular}
\caption{Separate participant-model usage ledgers. Cached input is already
included in input, and reasoning tokens, when reported, are part of output.
Neither category is added a second time. Pilot usage is in Appendix D.}
\label{tab:panel-usage}
\end{table}

Table~\ref{tab:panel-usage} reports archived receipt-derived counts, not
analyst usage or provider charges; these latter quantities were not
measured. Synthetic transport and privileged local computation are outside
the table. The 32-trial plan reserves 1,050,000 input plus 128,000 output
tokens for each request lacking usage, separately from measured totals.
Its final unknown-call count is zero, so no such reserve remains in that
panel's final accounting; this policy is not a provider billing cap.
The source-bound scalar projection is
\texttt{artifact/data/protocol-summary.json}. It provides provenance and
protocol identities without redistributing the full provider logs.


\paragraph{Exact requests and executable reader path.}
\texttt{artifact/protocol/} contains both condition instructions, the shared
wire schema, the four initial packets and all 80 request payloads, including
each trial's accumulated analysis records. Its checker renders every full
prompt and verifies the archived prompt and payload hashes. For example,
\begin{verbatim}
python3.12 -I -B -S artifact/protocol/check.py --show-call 1
\end{verbatim}
prints the first request. The historical collection entry point is supplied
as inspection-only text; provider transport credentials and full provider
responses are not redistributed. Configured profile names are not independent
attestation of server weights.

\texttt{artifact/menu\_replay.py} binds each public packet to its transfer
history and menu, then recomputes the 26 continuous and six route menu labels.
The original record comparison and this numerical replay are separate checks.
Continuous loss labels retain the open uniform-library and historical-runtime
premises in Appendix F; present-host agreement does not discharge them.

\section{Study inventory and historical boundaries}
\label{app:inventory}
The eight-trial pilot, 48-trial prospective transfer panel, 32-trial
common-menu panel and separate 48-assignment follow-up are four evidence sets
in the main paper. The two 48-slot studies are distinct. Their counts are
never pooled. The follow-up has 46 completed trials and two technical losses,
not 48 completed scientific trials. The two 1,600-draw development samples and the
5,000-draw main acquisition study appear only in the companion; the earlier
prediction criteria failed and are retained there. The acquisition result
is not an additional live LLM or crossing experiment.

Older archive labels describe different work: \emph{H4} denotes an unfinished
collection whose 84 remaining assignments stayed paused and whose crossing
status is unassessed due to missing evidence. The \emph{historical technical
null} denotes an earlier infrastructure-limited record, not a successful
zero-use scientific trial; its unresolved usage reservation stays outside
the pilot accounting. These records supply no extra evaluable trials here.
The archive's v1.8/v1.9 labels identify preserved predecessor research
versions, not additional studies or independent replications. Current paper
revision numbers track presentation and reader-access changes. This revision
additionally incorporates the already completed follow-up as a separate
evidence set; predecessor scientific inputs, denominators and the primary
acquisition procedure remain unchanged.

\section{Analysis-interface errors and consumed slots}
\label{app:interface-audit}
The historical shared instruction describes \texttt{history.states} and
\texttt{history.actions}; the evaluator actually supplies a Python dictionary.
Five analysis calls failed on \texttt{history.states}. Two other calls failed
because \texttt{next} is absent from the restricted built-ins; the historical
instruction refers to the existing interface without listing those built-ins.
These are observed interface-use errors, not missing final decisions. The
seven affected trials are listed below. Each made three analysis attempts:
one failed and two succeeded. Errors consumed an attempt and were visible
to the next request. They did not trigger replacement trials or extra slots.

\begin{table}[htbp]\centering\small
\begin{tabular}{lllll}\toprule
History & Profile & Condition & Repeat & Error\\\midrule
probe-01 & gpt-6-astra & Mandatory & 1 & Dictionary attribute\\
probe-01 & gpt-6-astra & Optional & 2 & Dictionary attribute\\
probe-01 & gpt-6.1-sol & Mandatory & 2 & Dictionary attribute\\
probe-01 & gpt-6.1-sol & Optional & 1 & Unavailable \texttt{next}\\
probe-01 & gpt-6.1-sol & Optional & 2 & Unavailable \texttt{next}\\
probe-02 & gpt-6-astra & Optional & 1 & Dictionary attribute\\
probe-02 & gpt-6.1-sol & Optional & 2 & Dictionary attribute\\\bottomrule
\end{tabular}
\caption{All seven affected common-menu trials. All are continuous-task
trials; both profiles requested max reasoning effort. The first row is the
sole rule-sensitive mandatory record. These are descriptive associations,
not estimates of the errors' effects on actions or records.}
\end{table}

The 32 common-menu trials produced 48 analyses: 41 successful and seven
erroneous. In the separate transfer panel, all 24 route trials made zero
analysis attempts, while the 24 continuous trials made 72 attempts: 67
successful and five erroneous. These denominators are not pooled. The
frozen adapter still supports only one transfer record and no complete pair.

\texttt{artifact/interface\_audit.py} reconstructs all 32 trial rows from the
last cumulative analysis list in each trial's archived requests, checks every
earlier list as a prefix, and reconciles attempt counts with the final trial
result and remaining-slot field. It also accounts for the 48 transfer trials.
No archived submitted program is executed. This procedure avoids counting
the same error again each time it appears in a later prompt.

The reconstructed tables and their crossing identities remain checkable.
However, the observed scarcity and row structure occurred under this
imperfect interface. A corrected description could change analysis use,
record structure or actions; no such behavioral effect has been measured.
Local checks of dictionary access and built-in availability cannot answer
that causal question. A prospective comparison must preserve old records
and compare the two interface descriptions under a separately specified
allocation and budget. Mandatory reporting and external execution would
still be a combined intervention in either interface arm.

\paragraph{Pilot sensitivity reader boundary.}
\texttt{artifact/pilot\_sensitivity.py} regenerates all twelve single-candidate
plans and their five selection metrics from the disclosed component inputs,
then reconstructs the five selected indices and native-action compatibility.
The corresponding P/L outcomes are linked to the archived sensitivity grades;
this added check does not rerun those physical graders. It therefore closes
the plan-to-selector link while retaining the numerical and historical-runtime
limits of the grade evidence.


\section{Follow-up contract, accounting and verification boundary}
\label{app:followup}
The follow-up profiles are gpt-6-astra/max/Fast and gpt-6.1-sol/max/Fast.
Each has 24 original assignments, eight per arm and eight matched triplets.
The same four historical contexts are used, with two repeats per profile;
neither repeats nor individual requests are new independent task samples.
A and B execute a legal proposal. A may omit its record; an invalid optional
record does not invalidate an otherwise legal proposal. B and C require a
valid record and legal proposal; a protocol-invalid final produces no action.
C executes the declared record selector and retains the proposal. The common
schema checks case/menu identity, dimensions, finite values and legal actions.
Single rows, duplicates and same-action rules are allowed. The shared
analysis interface uses dictionary access and states its available operations,
addressing the earlier interface mismatch; this is not a randomized old/new
interface comparison, so its causal effect is unmeasured.

Each trial permits at most three optional analyses and one final decision,
four model requests, 900 seconds per request and 3,780 seconds per trial.
Errors consume a slot; model responses are not repaired and failed
assignments are not replaced. The closed results record 36 analyses for
Astra and 36 for Sol, with zero recorded analysis errors in each. These counts
do not turn the historical interface comparison into a controlled experiment.
Decisions were bound before privileged label lookup. Follow-up labels reuse
the same identity-bound finite menu census, retaining the continuous numerical
and historical-runtime premises of the earlier paper. No new physics or
live crossed model decision was generated for this analysis.

\paragraph{Three execution segments and missingness.}
The original 48-position allocation was retained across three separately
authorized execution segments. Positions 1--12 completed; 13 and 14 ended
with HTTP 429 technical losses in the first and second segments. The third
segment completed positions 15--48 without repeating either failed position.
Both technical losses are Astra, in A and C of the same matched triplet.
Thus Astra has 22 completed, two technical and zero protocol losses; Sol has
24 completed and zero technical or protocol losses. All 48 assignments started
and have result records; none remains active or unstarted. Completion of the
34-position last segment does not mean all 48 assignments scientifically
completed. Missing pairs are excluded only from observed-pair descriptions,
with their denominators visible; no missing-at-random claim is made.
The final segment's 88 request receipts contain no terminal 429 diagnostic.
That observation does not identify the earlier limiter or establish a
permanent resolution; credit exhaustion is not inferred.

\begin{table}[htbp]
\centering\small
\begin{tabular}{lrrrrr}\toprule
Profile & Requests & Known/unknown & Input & Cached input & Output\\\midrule
Astra & 60 & 58/2 & 1,428,745 & 311,168 & 305,022\\
Sol & 60 & 60/0 & 1,418,971 & 263,424 & 294,997\\\bottomrule
\end{tabular}
\caption{Follow-up accounting only. Token counts cover known usage.
Cached input is included in input; reasoning output is included in output.
Neither subset is added again. Requests are not independent trials.}
\label{tab:followup-usage}
\end{table}

Using the frozen tariffs, measured credit equivalents are 1,336.90190 for
Astra and 264.37032 for Sol. Astra's two unknown requests retain a separate
1,690 credit-equivalent reserve, based on 1,050,000 input plus 128,000 output
tokens per request. These are local accounting quantities, not a provider
invoice or provider token ceiling. The original 7,000 local admission bound
was not increased. Measured usage and unknown reserves remain separate;
cache splits are measured for known receipts, while any cold-price treatment
of an unknown split would be an assumption, not measured zero cache.

\paragraph{Schema denominators.}
The retained field \texttt{partial\_triplets} counts planned blocks missing
an A, B or C assignment; its zero value does not assert scientific completion.
The reader separately reports fully completed triplets (Astra 7, Sol 8),
scientifically incomplete triplets (Astra 1, Sol 0), and triplets containing
technical losses (Astra 1, Sol 0). The retained field
\texttt{cache\_split\_unknown\_requests} counts unknown cache splits only
among requests with known token usage: its assessed denominators are
58 for Astra and 60 for Sol. Astra's two usage-unknown requests are excluded
from that count and retain their full separate reserves. Thus zero unknown
cache splits is neither zero cached input nor known usage for every request.
The machine-readable \texttt{artifact/followup/SCHEMA.json} defines these
populations and exclusions explicitly.

\paragraph{Reader inputs and remaining scope.}
The package contains 46 exact minimal finals, including 31 original numeric
tables and declared rules, plus the common/A/B/C instructions and wire schema.
The reader checks trial, profile, block, case and menu identities; reparses
decimal tokens as exact rationals; recomputes rule selections and row diagnostics;
and reconciles decisions, commit hashes and aggregates. It checks A/B against
the proposals and C against the record selector. All 46 completed labels are
re-read from the bound \texttt{artifact/data/crossing.json}; none changed.
Section~\ref{sec:reader-access} gives the source paths and executable entry point. No participant
analysis or task physics is run by this source-table check.

Full archives remain local. Hashes do not establish independent laboratory
replication or immutable provider-model/tier identity. The proposed population,
live query-selection and representation-transfer studies were not run;
the finite-acquisition companion retains its separate, unchanged sample and claims.


\subsection{History-level outcome accounting}
\label{app:followup-history}
Table~\ref{tab:followup-histories} retains the four purposeful histories and
their two repeats within each profile. It describes these assignments only;
neither the cells nor their repeated contexts are independent new worlds.
The reader derives this table and the main-text family, transition and
rule-sensitivity tables from the bound minimal finals, allocation and menu
census. Technical assignments retain no scientific label. All completed
follow-up labels are P or L; there are no U or protocol outcomes.
% Generated from identity-bound existing follow-up sources.
\begin{table}[htbp]
\centering\small
\begin{tabular}{llllll}\toprule
Profile & History & Family & A: P/L/T & B: P/L/T & C: P/L/T\\\midrule
Astra & probe-01 & Continuous & 1/1/0 & 2/0/0 & 1/1/0\\
Astra & probe-02 & Continuous & 2/0/0 & 2/0/0 & 1/1/0\\
Astra & probe-03 & Route & 2/0/0 & 2/0/0 & 2/0/0\\
Astra & probe-04 & Route & 1/0/1 & 2/0/0 & 1/0/1\\
Sol & probe-01 & Continuous & 1/1/0 & 0/2/0 & 2/0/0\\
Sol & probe-02 & Continuous & 2/0/0 & 2/0/0 & 1/1/0\\
Sol & probe-03 & Route & 2/0/0 & 2/0/0 & 2/0/0\\
Sol & probe-04 & Route & 2/0/0 & 2/0/0 & 2/0/0\\
\bottomrule
\end{tabular}
\caption{All follow-up history/arm outcomes; T denotes technical loss with no scientific label. Every cell has two planned assignments. All completed labels are P or L. Repeated contexts are not independent draws of new worlds; P is continuation feasibility, not measured full-task completion.}
\label{tab:followup-histories}
\end{table}

